\documentclass[12pt]{article}
\usepackage[T1]{fontenc}
\usepackage[utf8]{inputenc}
\usepackage{lmodern}
\usepackage{textcomp}
\usepackage{enumitem}
\usepackage{geometry}
\geometry{margin=1in}
\begin{document}
\begin{center}
Answer Key - Political Science - Graduate Level Assessment 3 \
\small (Answers are ordered exactly as in the question paper.)
\end{center}

\section*{Multiple Choice}
\begin{enumerate}[label=\arabic*.]
\item \textbf{Answer:} A
\\ \textit{Options:} A) Realism; B) Idealism; C) Liberalism; D) None of the above
\\ \textit{Note:} Realism posits that states are the central players in international relations, emphasizing their pursuit of power and national interests in an anarchic global system.
\item \textbf{Answer:} C
\\ \textit{Options:} A) terrorism policy.; B) economic policy.; C) foreign policy.; D) international policy.
\\ \textit{Note:} Foreign policy encompasses the strategies and decisions made by a nation to manage its interactions and relationships with other countries, making it the appropriate term for the realm of policy decisions that focus on international relations.
\item \textbf{Answer:} C
\\ \textit{Options:} A) 1987; B) 1997; C) 1994; D) 2004
\\ \textit{Note:} The Human Development Report, which introduced the Human Development Index and emphasized a broader understanding of development beyond economic growth, was first published by the United Nations Development Programme in 1994.
\item \textbf{Answer:} B
\\ \textit{Options:} A) Sovereignty.; B) Identity.; C) All of these options.; D) Hegemonic ethnicity.
\\ \textit{Note:} Identity is essential to societal security because it fosters a sense of belonging and cohesion within a community, which is crucial for maintaining stability and mitigating conflicts arising from cultural, ethnic, or ideological differences.
\item \textbf{Answer:} B
\\ \textit{Options:} A) "mutual assured destruction."; B) "brinkmanship."; C) "realism."; D) "not in my backyard."
\\ \textit{Note:} The correct answer is "brinkmanship" because John Foster Dulles articulated this policy as a strategy of pushing a dangerous situation to the brink of disaster to achieve favorable outcomes in the Cold War, particularly in relation to the Soviet Union.
\end{enumerate}

\section*{True / False}
\begin{enumerate}[label=\arabic*.]
\setcounter{enumi}{5}
\item \textbf{Answer:} False
\\ \textit{Options:} A) True; B) False
\\ \textit{Note:} Liberal Internationalists advocate for active U.S. engagement in international organizations and alliances to promote cooperation, democracy, and peace, rather than advocating for withdrawal and independence from the global arena.
\item \textbf{Answer:} True
\\ \textit{Options:} A) True; B) False
\\ \textit{Note:} Motivational realism posits that states act primarily out of self-interest, leading them to pursue policies aimed at accumulating wealth and promoting their ideological frameworks as a means of enhancing their power and influence on the global stage.
\item \textbf{Answer:} False
\\ \textit{Options:} A) True; B) False
\\ \textit{Note:} The statement is false because while NATO plays a crucial role in military and security matters, other organizations such as the United Nations are more significant in terms of comprehensive global governance, addressing a wider range of issues including human rights, development, and international law.
\item \textbf{Answer:} True
\\ \textit{Options:} A) True; B) False
\\ \textit{Note:} Inefficient balancing or buckpassing by states can create power vacuums and instability, encouraging competition among states to assert dominance or influence in the international system.
\item \textbf{Answer:} True
\\ \textit{Options:} A) True; B) False
\\ \textit{Note:} Realists assert that the absence of a central authority in the international system leads to anarchy, which fundamentally differs from domestic systems where governments enforce order and laws.
\end{enumerate}

\section*{Long Form Response}
\begin{enumerate}[label=\arabic*.]
\setcounter{enumi}{10}
\item \textbf{Answer:} Braun and Chyba (2004) argue that indigenous nuclear programs pose significant challenges to the effectiveness of proliferation safeguards, as these programs often operate outside established international frameworks and oversight mechanisms. The authors suggest that countries developing their own nuclear capabilities may prioritize national security over compliance with non-proliferation treaties, thereby undermining the very safeguards designed to prevent the spread of nuclear weapons. Evidence supporting this perspective includes case studies of nations that have pursued nuclear development independently, demonstrating a tendency to evade scrutiny and manipulate regulatory frameworks. Furthermore, the authors highlight the technological advancements that enable states to conceal aspects of their nuclear programs, further complicating verification efforts. Overall, the interplay between national interests and the limitations of existing safeguards illustrates the complexities introduced by indigenous nuclear initiatives in the global non-proliferation landscape.
\\ \textit{Note:} correct because it accurately captures Braun and Chyba's assertion that indigenous nuclear programs can weaken proliferation safeguards by prioritizing national security interests over international compliance, thereby highlighting the challenges these programs pose to global non-proliferation efforts.
\item \textbf{Answer:} World War I catalyzed significant economic transformations, notably shifting economic power from Europe to the United States. As European nations engaged in extensive war spending, they increasingly turned to American banks and businesses for loans and supplies, effectively establishing the U.S. as a creditor nation. This influx of capital and the demands of wartime production stimulated the American economy, allowing it to emerge as a leading financial center. The implications of this shift included a redistribution of global economic influence, with the U.S. gaining unprecedented leverage in international financial matters and setting the stage for its dominant role in the interwar period and beyond. Thus, the war not only altered immediate economic conditions but also redefined the global economic landscape for decades to come.
\\ \textit{Note:} correct because it clearly outlines how World War I led to a transfer of economic dominance from Europe to the United States through increased American lending and wartime production, establishing the U.S. as a creditor nation and a leading financial center.
\end{enumerate}

\vfill
\noindent\textit{End of Answer Key}
\end{document}